\section{Excluding finite abelian characters}\label{sec:characters}

Throughout this section, $L/k$ is quadratic, $D$ is a central division
algebra over $L$ of odd degree $n\geq3$, and $*$ is a unitary involution.
We retain $U=\U(D,*)$ and $S=\SU(D,*)$ from
Section~\ref{sec:colors}. Our aim is the following assertion, which
eliminates the abelian alternative of
Proposition~\ref{prop:color-dichotomy}.

\begin{theorem}\label{thm:no-characters}
Every homomorphism from $U(k)$ to a finite abelian group is trivial on
$S(k)$.
\end{theorem}

Fix such a homomorphism $\chi:U(k)\to B$, and write $B$ additively.
The proof first converts $\chi$ into a difference function on $D$.
Exact elementary transformations make this difference independent of its
arguments. A real-signature argument then converts that algebraic
identity into continuity on $S(k)$, after which local commutators finish
the proof. No continuity of $\chi$ is assumed.

\subsection{A noncommutative difference identity}

Let $\mathcal D$ be the set of nonzero $b\in D$ for which
$b+b^*=a^*a$ for some $a\in D$. Define
\begin{equation}\label{eq:chars-f}
 m_b=1-ab^{-1}a^*,\qquad f(b)=\chi(m_b).
\end{equation}
These definitions also allow $a=0$, in which case $b$ is skew and
$m_b=1$. Multiplication using $a^*a=b+b^*$ gives
$m_bm_b^*=1$. If $a'{}^*a'=a^*a\ne0$, then $a'=va$ with
$v\in U(k)$, and the two resulting elements $m_b$ are conjugate by
$v$. Thus $f$ is well-defined. Moreover,
\begin{equation}\label{eq:chars-congruence}
 f(x^*bx)=f(b)\qquad(x\in D^\times),
\end{equation}
because $ax$ can be used in the definition on the left.

For the moment fix $d\in D^\times$ such that the reduced spectra of
$d$ and $d^*$ are disjoint. For $(a,c)\in D^2\setminus\{(0,0)\}$,
there is a unique $e\in D$ satisfying
\begin{equation}\label{eq:chars-sylvester}
 e+e^*=c^*c,\qquad ed+d^*e^*=a^*a.
\end{equation}
Indeed the equation
\[
 ed-d^*e=a^*a-d^*c^*c
\]
has a unique solution: after splitting the algebra its linear part is
the Sylvester operator, whose eigenvalues are the differences of the
two disjoint spectra. Taking adjoints shows that $c^*c-e^*$ is another
solution of this same equation, giving the first identity in
Equation~\ref{eq:chars-sylvester} and then the second. The solution is
nonzero, for $e=0$ would imply $a=c=0$. Consequently
\[
 F_d(a,c)=f(ed)-f(e)
\]
is defined on all nonzero pairs.

\begin{lemma}\label{lem:chars-pair-invariance}
The function $F_d$ is invariant under independent left multiplications
of its two arguments by elements of $U(k)$, and under
\begin{equation}\label{eq:chars-pair-move}
 (a,c)\longmapsto(a-cd,a-c).
\end{equation}
\end{lemma}

\begin{proof}
The independent unitary multiplications leave
Equation~\ref{eq:chars-sylvester} unchanged. For the second assertion put
$b=ed$ and $z=b^*+e-a^*c$. Expansion, with the order of all factors
unchanged, gives
\[
 z+z^*=(a-c)^*(a-c),\qquad
 zd+d^*z^*=(a-cd)^*(a-cd).
\]
The transformation in Equation~\ref{eq:chars-pair-move} is invertible:
its kernel would give $a=c$ and $c(1-d)=0$, whereas $d\ne1$ by spectral
disjointness. In particular $z\ne0$.

Set $A'=ce^{-1}-ab^{-1}$ and $q=b^{-*}ze^{-1}$, where
$b^{-*}=(b^*)^{-1}$. A direct expansion gives $A'{}^*A'=q+q^*$.
Also
\[
 q=(b^{-1})^*(zd)b^{-1},
\]
so Equation~\ref{eq:chars-congruence} identifies $f(q)$ with $f(zd)$.
For clarity, the multiplicative identity needed here is
\begin{equation}\label{eq:chars-three-unitaries}
 1-(A'e)z^{-1}(b^*A'{}^*)=m_bm_zm_e^*,
\end{equation}
where the defining vectors for $m_b,m_z,m_e$ are respectively
$a,c-a,c$. It follows from the two identities
\begin{align*}
 A'e&=-ab^{-1}z+m_b(c-a),\\
 b^*A'{}^*&=ze^{-*}c^*+(c-a)^*m_e^*.
\end{align*}
To check the remaining terms after multiplying these identities, use
$m_ba=-ab^{-1}b^*$ and $c^*c=e+e^*$. The part of the left side of
Equation~\ref{eq:chars-three-unitaries} not containing
$-(m_b(c-a))z^{-1}((c-a)^*m_e^*)$ is then $m_bm_e^*$.
The left side of Equation~\ref{eq:chars-three-unitaries} is precisely
the unitary element defining $f(q)$, because $q^{-1}=ez^{-1}b^*$.
Applying $\chi$ proves
\[
 f(zd)=f(b)+f(z)-f(e),
\]
which is the required invariance. These computations take place in
$D$; no commutativity of their entries is used.
\end{proof}

\subsection{Producing exact elementary transformations}

Suppose now that $d$ lies in a $*$-stable maximal subfield $E\subset D$
which is Galois over $L$. Write $\bar z=z^*$ for $z\in E$, and suppose
that the $n$ conjugates $t_1,\ldots,t_n$ of $d$ are distinct and disjoint
from $\bar t_1,\ldots,\bar t_n$. Skolem--Noether gives a decomposition
\[
 D=\bigoplus_{j=1}^n E w_j,\qquad w_jd=t_jw_j.
\]
Let $E^1=\{z\in E^\times:z\bar z=1\}$. The pair move in
Equation~\ref{eq:chars-pair-move} and independent left multiplications
by elements of $E^1$ preserve the $n$ two-dimensional $E$-blocks in
$D\oplus D$.

In the invariance group take the subgroup $H$ generated by diagonal
$E^1$-multiplications on either coordinate and by the simultaneous matrices
\begin{equation}\label{eq:chars-boosts}
 B_h(t_j)=\begin{pmatrix}1&t_j\bar h\\h&1\end{pmatrix}
 \qquad(h\in E^1).
\end{equation}
Indeed changing the sign of the second output in
Equation~\ref{eq:chars-pair-move} gives $B_{-1}$, and conjugating by
diagonal phases gives all $B_h$. Each determinant is $1-t_j\ne0$.

\begin{lemma}\label{lem:chars-pure-shears}
For each block $i$, $H$ contains a nonidentity upper elementary
unipotent supported only on block $i$, and a nonidentity lower
elementary unipotent supported only on block $i$.
\end{lemma}

\begin{proof}
We first separate one block from the others. On a block with parameter
$t$, the product $B_h(t)B_{hz}(t)$ sends the first coordinate axis to
the line of slope
\[
 h\frac{1+z}{1+tz}\qquad(z,h\in E^1).
\]
At $t=t_j$ the norm of the slope with $h=1$ is
\begin{equation}\label{eq:chars-slope-norm}
 \ell_j(z)=\frac{(1+z)^2}{(1+t_jz)(z+\bar t_j)}.
\end{equation}
This is a degree-two rational function. Apart from finitely many $z$,
the other solution $z'$ of $\ell_j(z')=\ell_j(z)$ is distinct from
$z$ and lies in $E$. Explicitly, with $a=t_j$ and $b=\bar t_j$,
\[
 z'=\frac{(a-b)z+1+ab-2b}{(1+ab-2a)z+b-a}.
\]
It belongs to $E^1$: conjugate reciprocation
preserves the two roots and fixes the first, so fixes the second.
If $s_t(z)=(1+z)/(1+tz)$, choose
$h'=s_{t_j}(z)/s_{t_j}(z')\in E^1$. Then
\[
 g_j(t)=\bigl(B_{h'}(t)B_{h'z'}(t)\bigr)^{-1}B_1(t)B_z(t)
\]
has lower-left entry zero at $t_j$.

The lower-left numerator of $g_j(t)$ is
\[
 (1+tz')(1+z)-h'(1+z')(1+tz).
\]
This is a nonzero linear polynomial in $t$: the corresponding two
fractional slope functions have distinct poles and nonzero numerators.
Its unique root is $t_j$. Thus the
lower-left entry of $g_j(t)$ is nonzero at every other $t_\ell$ and at
every $\bar t_\ell$. Matrices of the form
\[
 \begin{pmatrix}\alpha(t)&t\overline{\beta}(t)\\
                 \beta(t)&\overline{\alpha}(t)\end{pmatrix},
\]
where the bar acts on coefficients and fixes the indeterminate $t$,
are closed under multiplication and inversion. The boosts have this
form. Consequently the upper-right entry of $g_j(t)$ vanishes exactly
when its lower-left entry vanishes at $\bar t$. It is therefore
nonzero at all $t_\ell$, including $t_j$.

We may also arrange that all diagonal entries at every $t_\ell$ are
nonzero. Here is a check that the excluded conditions are proper.
At $z=-1$, Equation~\ref{eq:chars-slope-norm} has a double zero and its
degree-two interchange sends $z$ to $z'$ with derivative $-1$.
Thus $h'\to-1$ as $z\to-1$, and both boost products tend to
$(1-t)I$. Hence $g_j(t)\to I$. These are algebraic limits on the
norm-one parameter curve over the fixed field of bar; its rational
points $E^1$ are infinite, by Hilbert 90, so avoid the finitely many
proper exceptional conditions.
We have obtained a triangular, nondiagonal matrix at $j$, and matrices
with all four entries nonzero at the other blocks. Repeating the
construction with $\bar t_j$ in place of $t_j$ gives a lower triangular,
nondiagonal matrix at $j$.

Fix a block $i$. For each $j\ne i$, let $H^{(j)}$ be the subgroup
generated by $g_j$ and the diagonal phases. Its projection at $j$ is
upper triangular. Its projection at $i$ is projectively Zariski dense
in $\PGL_2$ over an algebraic closure of $E$: it contains a
Zariski-dense subgroup of the diagonal
torus and a matrix with all entries nonzero. A proper algebraic subgroup
containing the diagonal torus lies in a Borel subgroup or its torus
normalizer, neither of which contains that matrix. Since the second
derived group of a triangular group is trivial,
$(H^{(j)})''$ acts trivially at $j$, but remains projectively Zariski
dense at $i$. The latter assertion follows by taking the Zariski
closures of commutators in the perfect algebraic group $\PGL_2$.

Importantly, use the \emph{full} kernels
\[
 K_j=\ker\bigl(H\longrightarrow\GL_2(E)
                         \text{ on block }j\bigr).
\]
Each $K_j$ is normal in $H$ and contains $(H^{(j)})''$, so its image
at $i$ is projectively dense. A commutator of two such kernels is
trivial on the union of their excluded blocks and remains projectively
dense at $i$: the commutators of two dense images are dense in the
derived algebraic group $\PGL_2$. Iterating gives a projectively dense
image at $i$ of
\[
 K_{\ne i}=\bigcap_{j\ne i}K_j.
\]
Write $H_i$ for the projection of $H$ at $i$, and $P_i$ for the image
of $K_{\ne i}$ there. We have proved that $P_i\triangleleft H_i$ and
that $P_i$ is projectively Zariski dense.

We next pass from this density to actual unipotents. Commuting the
upper and lower triangular, nondiagonal elements already constructed
with diagonal phases shows that $H_i$ contains upper and lower
elementary unipotents with some nonzero coefficients. Phase conjugation
and addition show that their coefficients contain scalar multiples of
the additive group $\mathbb Z E^1$. Its rational span is all of $E$.
Indeed that span is a finite-dimensional subring of a field, hence a
field; the Cayley ratios $(1+x)/(1-x)$ for $\bar x=-x$ show that it
contains the whole skew subspace, which generates $E$. It follows that
the allowable coefficients for both upper and lower unipotents contain
a common nonzero integral ideal $I\subset\mathcal O_E$.

Choose $p\in P_i$ with $p\infty=r\ne\infty$ on $\mathbb P^1(E)$.
Choose an infinite-order unit $u\in\mathcal O_E^\times$ sufficiently
close to $1$ modulo an integral ideal that
\[
 u=1+st\quad(s,t\in I),\qquad (1-u^2)r\in I.
\]
Such units exist by the unit theorem: $[E:\mathbb Q]\ge6$, and a
congruence subgroup of the unit group has finite index. More explicitly,
fix $0\ne s\in I$, require $u-1\in sI$, and impose the additional
congruence clearing the denominator of $r$. With
$U(x)=\left(\begin{smallmatrix}1&x\\0&1\end{smallmatrix}\right)$
and $L(x)=\left(\begin{smallmatrix}1&0\\x&1\end{smallmatrix}\right)$,
the exact identity
\[
 U(s)L(t)U(-s/u)L(-tu)=\diag(u,u^{-1})
\]
shows that $H_i$ contains the affine transformation
\[
 h:x\longmapsto u^2x+(1-u^2)r.
\]
It fixes $r$ and $\infty$. Normality gives
$q=h^{-1}p^{-1}hp\in P_i$. This element fixes $\infty$ and has affine
slope $u^{-4}$: in a local coordinate at $\infty$ the derivative of
$h$ is $u^{-2}$, whereas at $r$ it is $u^2$.
Thus, for any nonzero $x\in I$,
\[
 qU(x)q^{-1}U(-x)=U((u^{-4}-1)x)
\]
is an actual nonidentity upper unipotent in $P_i$. Scalar factors in
a representative of $q$ cancel in this identity. By the definition of
$P_i$ it lifts to an element of $H$ which is the identity on every
other block. Interchanging the axes proves the lower assertion.
\end{proof}

To turn a single nonzero shear coefficient into an arbitrary additive
change in $D$, we need the following exact statement, not merely local
density.

\begin{lemma}\label{lem:chars-additive-span}
The additive subgroup generated by $U(k)$ is all of $D$:
\[
 \mathbb Z U(k)=D.
\]
\end{lemma}

\begin{proof}
Let $M=\mathbb Z U(k)$. It is a subring, since $U(k)$ is a group.
For a rational prime $p$, let $C_p$ be its additive closure in
$D_p=D\otimes_{\mathbb Q}\mathbb Q_p$. Weak approximation for $U$
puts the product of its local groups above $p$ inside $C_p$. The largest
$\mathbb Q_p$-vector subspace of $C_p$ is
\[
 V_p=\bigcap_{j\ge0}p^jC_p.
\]
It is closed; here closed additivity implies stability under
$\mathbb Z_p$. Left and right multiplication by every local unitary
element preserve $V_p$. Differentiating these actions shows stability
under multiplication by every skew element, separately in each local
component. The skew elements generate the local algebra: an invertible
central skew scalar and its inverse, together with skew multiples of
Hermitian elements, generate all Hermitian and skew elements. Hence
$V_p$ is a two-sided ideal of $D_p$.

Every simple factor occurs in this ideal. At a nonsplit finite place of
$k$, the algebra over $L$ splits and the unitary group is isotropic
because $n\ge3$. Choose an unbounded sequence in that local unitary
group, taking the identity in the other local components. Multiplication
by suitable positive powers $p^{r_\nu}$, with $r_\nu\to\infty$, gives
a subsequence tending to a nonzero vector in this factor and zero in
the others. For each fixed $j$, the same sequence multiplied by
$p^{r_\nu-j}$ lies in $C_p$ eventually, so its limit is divisible by
$p^j$ in $C_p$. The nonzero limit belongs to $V_p$.
At a split place the two factors are interchanged by $*$; the unitary
group contains reciprocal central scalars. Taking $(p^{-r},p^r)$ and
scaling by $p^r$ gives a nonzero limit in one factor alone; reversing the
two scalars gives the other. Thus the ideal $V_p$ contains every simple
factor, proving $C_p=D_p$.

This local density implies that the rational span of $M$ is $D$, so
$M$ contains a full $\mathbb Z$-lattice $\Lambda$ in $D$. Fix
$x=p^{-j}\lambda$, with $\lambda\in\Lambda$. Choose $a\in M$ with
$a-x\in\Lambda\otimes\mathbb Z_p$. There is a positive integer $m$
prime to $p$, congruent to $1$ modulo $p^j$, which clears all denominators
of $a$ away from $p$ relative to $\Lambda$. Then
$ma-mx\in\Lambda$ at every rational prime, and hence globally, while
$(m-1)x\in\Lambda$. Therefore $x\in M$. Inverting all rational primes
in $\Lambda$ now gives $M=D$.
\end{proof}

\begin{proposition}\label{prop:chars-difference}
Let $E\subset D$ be a $*$-stable maximal subfield, Galois over $L$.
Let $d\in E^\times$ have $n$ distinct $L$-conjugates
$t_1,\ldots,t_n$ whose set is disjoint from
$\{\bar t_1,\ldots,\bar t_n\}$, where $\bar z=z^*$ on $E$.
Then, whenever $b,bd\in\mathcal D$,
\begin{equation}\label{eq:chars-difference}
 f(bd)-f(b)=\chi(\bar d/d).
\end{equation}
\end{proposition}

\begin{proof}
We show that the invariance group of $F_d$ is transitive on nonzero
pairs. If $c\ne0$, one of its $E$-block coordinates is nonzero. The
corresponding pure upper shear of Lemma~\ref{lem:chars-pure-shears}
changes $a$ by a fixed nonzero element $y\in D$, leaving $c$ unchanged.
Conjugating this transformation by left multiplication of the first
coordinate by $v\in U(k)$ changes $a$ by $vy$. Products and inverses
therefore permit addition of every element of $(\mathbb ZU(k))y=D$,
by Lemma~\ref{lem:chars-additive-span} and division. The lower shears
likewise permit arbitrary changes of $c$ when $a\ne0$. These operations
connect every nonzero pair to a pair with both entries nonzero, and then
connect any two such pairs. Thus $F_d$ is constant.

At $e=(d-\bar d)^{-1}$, Equation~\ref{eq:chars-sylvester} is solved
by $c=0$, $a=1$. The unitary elements defining $f(e)$ and $f(ed)$ are
respectively $1$ and $\bar d/d$. This evaluates the constant and proves
Equation~\ref{eq:chars-difference}.
\end{proof}

\subsection{Scalar normalization and a cyclic maximal field}

Let $\Omega$ be the nonsplit real places of $k$. At $v\in\Omega$
write the involution as $x^*=H_v^{-1}\bar x^{\,t}H_v$ on
$M_n(\mathbb C)$. For a nonzero Hermitian element $h\in D$, call its
signature admissible when the signature of $H_vh$ is that of $H_v$ up
to overall sign, for every $v\in\Omega$. This does not depend on the
matrix realization. Define $\mathcal D^\#$ to consist of the nonzero
skew elements and the $b\in D$ for which $h=b+b^*\ne0$ has admissible
signature. Division makes every such $h$ invertible.

\begin{lemma}\label{lem:chars-normalization}
If $h=h^*\ne0$ has admissible signature and $s=\Nrd(h)\in k^\times$,
there is $a\in D^\times$ with
\begin{equation}\label{eq:chars-normalized-norm}
 a^*a=sh,\qquad \Nrd(a)=s^{(n+1)/2}.
\end{equation}
Consequently
\[
 f^\#(b)=f(sb)\qquad(b\in\mathcal D^\#)
\]
is well-defined using any $s\in k^\times$ for which $sb\in\mathcal D$.
It is invariant under congruence by $D^\times$ and multiplication by
$k^\times$.
\end{lemma}

\begin{proof}
At a nonsplit finite place the algebra splits; classification of
Hermitian forms says that its dimension and determinant modulo norms
determine its isometry class. The determinant ratio for $sh$ is
$s^{n+1}$, which is a norm since $n+1$ is even. At a nonsplit real
place, $\operatorname{sign}(s)$ is positive if $H_vh$ has the signature
of $H_v$, and negative if it has the opposite signature, because $n$ is
odd. Thus $sH_vh$ has precisely the signature of $H_v$. At split finite
places the assertion reduces to surjectivity of the local reduced norm;
at split archimedean places the algebra is split, again by odd degree.
These facts give local solutions of $a^*a=sh$. Their reduced norms can
be made exactly $s^{(n+1)/2}$, since the determinant map of the local
full unitary group is onto its norm-one group. The variety in
Equation~\ref{eq:chars-normalized-norm} is a torsor under the simply
connected group $S$, and so has a rational point by its Hasse principle.
We use here the local Hermitian classification and the simply connected
Hasse principle in their number-field forms
\cite{KMRT,PRbook}.

If both $sb$ and $tb$ are in $\mathcal D$ and $b$ is not skew, taking
reduced norms of the two defining Hermitian equations shows that
$(t/s)^n$ is a norm from $L^\times$. The quotient
$k^\times/N_{L/k}(L^\times)$ has exponent two. Since $n$ is odd,
$t/s$ itself is a norm, say $t/s=\bar z z$. Central congruence by $z$
and Equation~\ref{eq:chars-congruence} give $f(tb)=f(sb)$. For skew
$b$ both values are zero. This proves independence. Congruence and
scalar invariance follow by applying the same argument to any
permissible normalization.
\end{proof}

\begin{lemma}\label{lem:chars-cyclic-field}
There is a $*$-stable maximal subfield $E\subset D$ which is cyclic of
degree $n$ over $L$. Its norm-one torus
$T=\ker(N_{E/L}:\Res_{E/L}\mathbb G_m\to\mathbb G_m)$ has weak
approximation over $L$.
\end{lemma}

\begin{proof}
By the odd-degree case of Grunwald--Wang, choose a cyclic extension
$F/k$ of degree $n$ having full local degree at each place below the
Brauer support of $D$ \cite{Neukirch}; see also
\cite[Chapter~VIII, Theorem~2.4 and Corollary~2.5]{MilneCFT}.
It is disjoint from $L$.
The local invariant criterion shows that $E_0=LF$ splits $D$; since
$[E_0:L]=n$, the embedding criterion for central simple algebras gives
an embedding $\alpha:E_0\hookrightarrow D$ \cite{Reiner}.
Let bar on $E_0$ fix $F$ and restrict to the nonidentity automorphism
of $L/k$. Skolem--Noether supplies $h\in D^\times$ with
\[
 \alpha(\bar z)^*=h\alpha(z)h^{-1}\qquad(z\in E_0).
\]
Replacing $z$ by $\bar z$ and taking adjoints gives
\[
 \alpha(z)=h^{-*}\alpha(\bar z)^*h^*.
\]
Thus $h^*$ also intertwines the two embeddings, and
$c=h^{-1}h^*$ lies in $\alpha(E_0)$. On this field the adjoint satisfies
$x^*=h\bar xh^{-1}$, so
\[
 \bar c=h^{-1}c^*h=h^{-*}h=c^{-1}.
\]
Hilbert 90 gives $t/\bar t=c$. The replacement
$h\mapsto h\alpha(t)$ preserves the intertwining identity and is
Hermitian, because
$(h\alpha(t))^*=h\alpha(\bar t)c=h\alpha(t)$.

At a real place in $\Omega$, the odd cyclic extension $F/k$ splits
completely. For the Hermitian form with matrix $H_vh$, the $n$
eigenlines of $\alpha(E_0)$ are orthogonal. Replacing $h$ by
$h\alpha(z)$, with $z\in F^\times$, scales their Hermitian coefficients
independently by the real embeddings of $z$. Weak approximation in $F$
therefore makes its signature admissible at every such place, even
equal to that of $H_v$ if desired. Apply
Lemma~\ref{lem:chars-normalization} to obtain $a^*a=sh$. The embedding
$z\mapsto a\alpha(z)a^{-1}$ is $*$-stable and induces bar, since
\[
 (a\alpha(\bar z)a^{-1})^*
 =(a^*)^{-1}h\alpha(z)h^{-1}a^*
 =a\alpha(z)a^{-1}.
\]
Its image is the required $E$. Finally, for a cyclic generator $\sigma$
of $\Gal(E/L)$, Hilbert 90 parametrizes $T(L)$ by
$y/\sigma(y)$ with $y\in E^\times$. The same parametrization holds at
each completion, and weak approximation in $E$ proves weak
approximation for $T$.
\end{proof}

Fix this field $E$ and an integer $m\ge1$ annihilating $B$. Call
$d\in T(L)^m$ admissible if its conjugates over $L$ are distinct and
disjoint from their bars. These conditions specify a nonempty $k$-Zariski
open subset of $\Res_{L/k}T$: after extending scalars, the two
determinant-one eigenvalue
lists can be chosen independently, avoiding equality within or between
the lists. In particular admissible elements exist, since the rational
points of $T$ are Zariski dense and the power map is dominant.
Let $\mathcal C$ be their $D^\times$-conjugates. This set is closed
under inversion.

\begin{corollary}\label{cor:chars-step-invariance}
If $t\in\mathcal C$ and $b,bt\in\mathcal D^\#$, then
\begin{equation}\label{eq:chars-step-invariance}
 f^\#(bt)=f^\#(b).
\end{equation}
\end{corollary}

\begin{proof}
Write $t=xdx^{-1}$. Congruence by $x$ gives
$x^*(bt)x=(x^*bx)d$. Choose separate scalar normalizations of these
two arguments. Their ratio only replaces $d$ in
Proposition~\ref{prop:chars-difference} by a nonzero $k$-multiple, which
preserves spectral disjointness and leaves $\bar d/d$ unchanged.
If $d=y^m$ with $y\in T(L)$, then
$\bar d/d=(\bar y/y)^m$, an $m$-th power in $E^1\subset U(k)$.
Thus its character is zero, and
Equation~\ref{eq:chars-step-invariance} follows.
\end{proof}

\subsection{Keeping products in the real signature domain}

The subgroup $J\subset\SL_1(D)(L)$ generated by $\mathcal C$ is
noncentral and normal. The established inner-type case of the
Margulis--Platonov theorem, applied over the number field $L$, gives an
open neighborhood $W$ of $1$ in
\[
 \prod_{w\in A_D}\SL_1(D)(L_w)
\]
whose rational inverse image is contained in $J$. Here $A_D$ is the set
of finite places of $L$ where $D\otimes_LL_w$ is division. This use of
the known inner case is independent of the outer case being proved
\cite[Theorem~2]{SegevSeitz2002}\cite[Sections~3.4--3.5]{PRsurvey}.
Although every element of $J$ is a word in the allowed steps, an
arbitrary word need not preserve the real signature domain at its
intermediate products. The next lemma supplies the required control.

\begin{lemma}\label{lem:chars-local-invariance}
Let $b\in\mathcal D^\#$ have nonzero Hermitian part. For every
$r\in\SL_1(D)(L)$ with $r_{A_D}\in W$ and $r_\Omega$ sufficiently
close to $1$, one has
\begin{equation}\label{eq:chars-local-invariance}
 f^\#(br)=f^\#(b).
\end{equation}
\end{lemma}

\begin{proof}
If $\Omega$ is empty, the signature condition is absent: every nonzero
element of $D$ is in $\mathcal D^\#$, either with nonzero Hermitian
part or skew. A word in $\mathcal C$ therefore proves the assertion
directly by Corollary~\ref{cor:chars-step-invariance}.

Suppose $\Omega\ne\varnothing$ and write
$G_\Omega=\prod_{v\in\Omega}\SL_n(\mathbb C)$. Let $\mathcal O$
be the open set of $g\in G_\Omega$ for which $bg$ has nondegenerate
Hermitian part with the required signatures. It contains $1$. Restrict
$r_\Omega$ to a sufficiently small neighborhood of $1$ that it can be
joined to $1$ by a path $\gamma:[0,1]\to\mathcal O$.
Since $r\in J$, write it as a word in elements of $\mathcal C$.
Join each factor of this word to $1$ in the corresponding conjugate of
$T$ at $\Omega$; these complex tori are connected. The product of the
factor paths is a path $p$ from $1$ to $r_\Omega$.

We explain why the based loop $p^{-1}\gamma$ is a finite product of
based loops in rational conjugates of $T$. The Lie algebras of finitely
many such conjugates span the real Lie algebra of $G_\Omega$.
Indeed the full conjugacy span is the whole Lie algebra by simplicity,
and rational conjugators in $D^\times$ are dense independently at the
finitely many places. The multiplication map from finitely many of
these tori consequently has a continuous local section near $1$, taking
$1$ to the tuple of identities. Applying this section pointwise factors
every based loop sufficiently uniformly close to $1$ into based torus
loops. The subgroup they generate in the based loop group is therefore
open. The based loop group is connected, because every
$\SL_n(\mathbb C)$ is simply connected; hence that open subgroup is
the whole based loop group. This proves the desired finite
factorization of $p^{-1}\gamma$.

Appending these factors to those of $p$ writes $\gamma$ pointwise as
a finite product of paths in rational conjugates of $T$, each starting
at $1$ and with rational prescribed endpoint. Every endpoint belongs
to the corresponding rational $m$-th-power group: the original word
has this property, and the appended loops end at $1$.

Take a fine mesh in $[0,1]$ with at least one interior point. At its
interior nodes approximate each torus factor by rational points in its
$m$-th-power group, leaving the endpoints unchanged. These rational
powers are dense at $\Omega$: use weak approximation for $T$ and
surjectivity of the power map on each complex torus. We can further
require that every consecutive ratio in every factor is admissible,
including the ratios next to the fixed endpoints. To see this precisely,
regard all interior nodes as independent variables in products of
$\Res_{L/k}T$. Every consecutive-ratio map is dominant; beside a fixed
endpoint it is simply translation or inverse translation of the free
variable. The complement of admissibility therefore pulls back to a
proper closed subset in each case. The finite union of these subsets
cannot fill the irreducible parameter variety. One may work in the
parameters before taking $m$-th powers, whose power map is dominant:
weak approximation then supplies rational parameters arbitrarily close
to the chosen local roots while avoiding the finite proper closed
exclusions. Thus both endpoint ratios and all interior ratios can be
made admissible simultaneously.

Pass from one mesh product to the next by changing one factor at a
time. Each change is right multiplication by a rational
$D^\times$-conjugate of an admissible increment. All intermediate
products remain in $\mathcal O$ if the mesh is sufficiently fine and
the approximations sufficiently close. To justify the uniform choice,
the finitely many factor paths are uniformly continuous, their images
are compact, and $\gamma([0,1])$ is a compact subset of the open set
$\mathcal O$; mixed products from two sufficiently nearby mesh nodes
are uniformly close to $\gamma$. At the rational intermediate points
the real nondegeneracy also guarantees nonzero Hermitian part in $D$.
Corollary~\ref{cor:chars-step-invariance} now applies to every move.
The initial and final products are exactly $1$ and $r$, so their
character differences give Equation~\ref{eq:chars-local-invariance}.
\end{proof}

\subsection{Continuity and the final vanishing}

\begin{proof}[Proof of Theorem~\ref{thm:no-characters}]
Choose $u\in U(k)$ with $\Nrd(u)\ne1$. Such an element can be chosen
central, since the norm-one torus of $L/k$ has infinitely many rational
points. Division implies that $1-u$ is invertible. Put
$b=(1-u)^{-1}$. Direct calculation gives
\begin{equation}\label{eq:chars-final-b}
 b+b^*=1,\qquad b^*=-ub,\qquad f^\#(b)=f(b)=\chi(u).
\end{equation}
If $u'\in U(k)$ has the same reduced norm as $u$, define
$b'=(1-u')^{-1}$; it satisfies the same identities. In particular
\[
 q=\frac{\Nrd(b')}{\Nrd(b)}\in k^\times.
\]
Indeed $\overline{\Nrd(b)}=-\Nrd(u)\Nrd(b)$ by oddness, and the same
identity for $b'$ makes the ratio fixed by $L/k$.

Suppose $u'$ is sufficiently close to $u$ at $A_D$ and at $\Omega$,
using both places of $L$ over every split place of $k$. Choose
$a\in D^\times$ with
\[
 \Nrd(a)=q^{(n-1)/2},
\]
and with $a$ close to $1$ at all these places. The global reduced norm
theorem gives one such element without the approximation conditions:
the only possible archimedean sign obstruction occurs at quaternionic
real factors, which cannot occur in odd degree. Locally near $1$ the
reduced norm is a submersion, so there are local solutions near $1$ for
$q$ near $1$. Weak approximation for the simply connected inner group
$\SL_1(D)$ then adjusts the global solution to these local solutions
\cite{PRbook,Reiner}.

Set $b''=q^{-1}a^*b'a$. Congruence and scalar invariance give
$f^\#(b'')=f^\#(b')$, while
\[
 \Nrd(b'')=q^{-n}q^{n-1}\Nrd(b')=\Nrd(b).
\]
Thus $r=b^{-1}b''\in\SL_1(D)(L)$ is close to $1$ at $A_D$ and
$\Omega$. Lemma~\ref{lem:chars-local-invariance} gives
$f^\#(b'')=f^\#(b)$, and hence $\chi(u')=\chi(u)$.
Taking $u'=us$ shows that $\chi|_{S(k)}$ is continuous for the topology
induced by these finitely many local factors.

The set $A_D$ consists precisely of the two places over each place in
$A$ from Section~\ref{sec:colors}. At nonsplit places the algebra
splits, and at a split place its local norm-one group is anisotropic
exactly when the component algebra is division. Weak approximation for
$S$ therefore extends this finite-valued continuous restriction uniquely
to a continuous homomorphism on
\[
 \prod_{v\in A}S(k_v)\ \times\ \prod_{v\in\Omega}S(k_v).
\]
For completeness, write $H$ for the dense rational subgroup and $K$ for
the kernel of its character. Continuity makes $K$ closed in $H$, so
$\overline K\cap H=K$. Finitely many cosets $h_iK$ cover $H$; the
finite closed union of $h_i\overline K$ therefore covers the ambient
local product. The subgroup $\overline K$ consequently has finite index
and is open, since its complement is a finite union of closed cosets.
It is normal by density and continuity of conjugation. The resulting
map $H/K\to G/\overline K$, where $G$ denotes the local product, is an
isomorphism: injectivity is the intersection identity, and surjectivity
is the coset cover. This gives the asserted continuous extension. Each real
factor is connected and so maps trivially to the finite group $B$.

Finally fix a split place $v\in A$, writing
$U(k_v)=\mathcal D_v^\times$ and $S(k_v)=\SL_1(\mathcal D_v)$.
Weak approximation in $U$ approximates any two local elements at $v$
by rational unitary elements, with the identity prescribed at the
other places of $A$ and $\Omega$. Their rational commutators belong to
$S(k)$ and have character zero. Continuity of the extended restriction
therefore kills every local commutator of $\mathcal D_v^\times$.
Lemma~\ref{lem:colors-local-abelianization} says that these commutators
generate a dense subgroup of $\SL_1(\mathcal D_v)$, so this factor also
maps trivially. All factors have been killed, proving
$\chi(S(k))=0$. This includes $A=\varnothing$, when only the connected
real factors occur, and includes the case where both products are
empty.
\end{proof}
